\documentclass[11pt]{article}

\usepackage[a4paper,margin=1in]{geometry}
\usepackage[T1]{fontenc}
\usepackage[utf8]{inputenc}
\usepackage{amsmath,amssymb,amsfonts,mathtools}
\usepackage{bm}
\usepackage{booktabs}
\usepackage{tabularx}
\usepackage{array}
\usepackage{algorithm}
\usepackage[noend]{algpseudocode}
\usepackage{hyperref}

\hypersetup{
    colorlinks=true,
    linkcolor=blue,
    citecolor=blue,
    urlcolor=blue
}

\newcommand{\Kset}{\mathcal{K}}
\newcommand{\Aset}{\mathcal{A}}
\newcommand{\Qset}{\mathcal{Q}}
\newcommand{\C}{\mathbb{C}}
\newcommand{\I}{\mathbf{I}}
\newcommand{\Hc}{\mathbf{H}}
\newcommand{\F}{\mathbf{F}}
\newcommand{\A}{\mathbf{A}}
\newcommand{\Sig}{\mathbf{\Sigma}}
\newcommand{\Smat}{\mathbf{S}}
\newcommand{\tr}{\operatorname{tr}}
\newcommand{\norm}[1]{\left\lVert #1 \right\rVert}
\newcommand{\floor}[1]{\left\lfloor #1 \right\rfloor}
\newcommand{\UL}{\mathrm{UL}}
\newcommand{\DL}{\mathrm{DL}}
\newcommand{\E}{\mathbb{E}}
\emergencystretch=1.5em
\allowdisplaybreaks

\title{Finite-sub-blocklength Latency Minimization for Uplink and Downlink Communications:\\Theoretical Formulation and Solution Methods}
\author{}
\date{\today}

\begin{document}

\maketitle

\begin{abstract}
This report presents the current uplink and downlink
work on finite-sub-blocklength latency minimization. It introduces a common
notation layer, defines the payload communication setting, derives the
finite-sub-blocklength rate expressions used in both links, formulates the uplink
and downlink latency-minimization problems, and then develops two solution
frameworks: the \emph{Training Only Method} and the \emph{Training + Testing
Method}. 
\end{abstract}

\section{Introduction}

Short-packet communication cannot be analyzed accurately through asymptotic
capacity alone. When the codeword length is finite, the achievable rate depends
not only on channel quality and transmit power, but also on sub-blocklength and the
required decoding reliability. For latency-sensitive multiuser communication,
this creates a coupled design problem: the number of transmitted bits, the
chosen sub-blocklength, and the precoder must be selected jointly.

The present work studies this problem in both uplink and downlink settings. In
the uplink, each user is handled through a user-wise sequence of transmission
blocks, and the main difficulty is that the number of required blocks is not
known a priori. In the downlink, all active users share the same transmission
epoch, so the main difficulty is inter-user interference and the resulting
coupling of feasible rates.

The report
\begin{enumerate}
    \item establish a common mathematical notation for uplink and downlink,
    \item define the payload communication setting that determines what each block is asked to deliver,
    \item derive the corresponding finite-sub-blocklength uplink and downlink formulations, and
    \item present the \emph{Training Only Method} and the \emph{Training + Testing Method} in algorithmic form.
\end{enumerate}


\section{System Model and Common Notation}

\subsection{Core Variables}

Table~\ref{tab:common-notation} lists the symbols that are shared throughout
the report.

\begin{table}[ht]
\centering
\small
\caption{Common notation}
\label{tab:common-notation}
\renewcommand{\arraystretch}{1.15}
\begin{tabularx}{\textwidth}{>{\raggedright\arraybackslash}p{0.27\textwidth} X}
\toprule
Symbol & Meaning \\
\midrule
$\Kset = \{1,2,\dots,K\}$ & Set of users. \\
$k \in \Kset$ & User index. \\
$\ell$ & Transmission-block index. In uplink it is per user; in downlink it is per Base Station. \\
$B_k$ & Total payload of user $k$ in bits. \\
$b_{k,\ell}^{\mathrm{req}}$ & Requested bits to send for user $k$ in block $\ell$. In the present report, $b_{k,\ell}^{\mathrm{req}} = B_{k,\mathrm{rem}}^{(\ell)}$. \\
$b_{k,\ell}$ & Bits actually sent to user $k$ in block $\ell$. \\
$B_{k,\mathrm{rem}}^{(\ell)}$ & Remaining unsent bits of user $k$ when block $\ell$ starts. \\
$n_{k,\ell}$ & sub-blocklength assigned to user $k$ in block $\ell$, measured in channel uses. \\
$T_k$ & Maximum admissible sub-blocklength for user $k$ within one coherence-limited transmission opportunity. \\
$f_{s,k}$ & Symbol rate of user $k$ in symbols per second. \\
$\tau_k$ & End-to-end latency of user $k$. \\
$\varepsilon_k$ & Target packet error probability of user $k$. \\
$P_k$ & Uplink power budget of user $k$. \\
$P_B$ & Total transmit-power budget of the downlink Base Station Precoder. \\
$d_k$ & Number of transmitted spatial streams of user $k$. \\
$\Aset_\ell$ & Active-user set in downlink block $\ell$. \\
$q$ & One query, i.e.\ one visited state presented to the learned model. \\
$\Qset$ & Collection of training queries used in the Training + Testing Method. \\
\bottomrule
\end{tabularx}
\end{table}

The latency of user $k$ is defined by
\begin{equation}
\tau_k
\frac{1}{f_{s,k}}
\sum_{\ell} n_{k,\ell},
\label{eq:latency-user}
\end{equation}
and the network objective is to minimize the aggregate latency
\begin{equation}
=
\sum_{k\in\Kset}\tau_k.
\label{eq:latency-sum}
\end{equation}

\subsection{Finite-sub-blocklength Template}

For both uplink and downlink, the achievable rate is described through the normal
approximation \cite{polyanskiy2010}. Once the specific  matrix
$\A_{k,\ell}^{(m)}$ has been defined, the corresponding capacity term is
\begin{equation}
C_{k,\ell}^{(m)}
=
\log_2\det\!\bigl(\I+\A_{k,\ell}^{(m)}\bigr),
\label{eq:common-C}
\end{equation}
and the channel-dispersion term is
\begin{equation}
V_{k,\ell}^{(m)}
=
(\log_2 e)^2
\tr\!\left[
\A_{k,\ell}^{(m)}
\bigl(\A_{k,\ell}^{(m)}+2\I\bigr)
\bigl(\A_{k,\ell}^{(m)}+\I\bigr)^{-2}
\right].
\label{eq:common-V}
\end{equation}

Therefore, for sub-blocklength $n_{k,\ell}$ and reliability target
$\varepsilon_k$, the finite-sub-blocklength rate is
\begin{equation}
R_{k,\ell}^{(m)}
=
C_{k,\ell}^{(m)}
-
\sqrt{\frac{V_{k,\ell}^{(m)}}{n_{k,\ell}}}\,
Q^{-1}(\varepsilon_k).
\label{eq:common-rate}
\end{equation}

If $b_{k,\ell}$ information bits are sent over $n_{k,\ell}$ channel uses,
feasibility requires
\begin{equation}
\frac{b_{k,\ell}}{n_{k,\ell}}
\le
R_{k,\ell}^{(m)}.
\label{eq:common-feasibility-rate}
\end{equation}
Multiplying both sides by $n_{k,\ell}$ yields
\begin{equation}
b_{k,\ell}
\le
n_{k,\ell}R_{k,\ell}^{(m)}.
\label{eq:common-feasibility-bits}
\end{equation}
Since the service variable is integer-valued, the largest feasible blockwise
service is
\begin{equation}
b_{k,\ell}^{\max,(m)}
=
\floor{n_{k,\ell}R_{k,\ell}^{(m)}}.
\label{eq:common-bmax}
\end{equation}

Equation~\eqref{eq:common-bmax} is central throughout the report: it converts a
continuous finite-sub-blocklength rate into an integer blockwise service limit.

\section{Communication Scenario}

Before separating uplink and downlink, it is useful to define the traffic
interpretation used throughout the report. This scenario determines what each
block is asked to transmit and how latency is accumulated across blocks.

\subsection{Payload Communication}

In payload communication, each user begins with a finite payload
$B_k$. The scheduler stops only when this payload is fully
drained. The number of bits to send in block $\ell$ is the current remainder:
\begin{equation}
b_{k,\ell}^{\mathrm{req}}=B_{k,\mathrm{rem}}^{(\ell)}.
\label{eq:payload-request}
\end{equation}
After the block is committed, the remainder is updated as
\begin{equation}
B_{k,\mathrm{rem}}^{(\ell+1)}=B_{k,\mathrm{rem}}^{(\ell)}-b_{k,\ell}.
\label{eq:payload-update}
\end{equation}
Thus, payload communication is a workload-draining problem: every decision is made
with respect to the residual bits that still have to be delivered.

\subsection{Uplink Payload}

For uplink payload communication, each user evolves independently according to its
own remainder recursion. User $k$ creates a new block only if its current block
cannot finish the entire remaining payload. Thus, the block count
$L_k$ is an output of the optimization.

\subsection{Downlink Payload}

For downlink payload communication, the active-user set changes with time:
\begin{equation}
\Aset_\ell = \{k\in\Kset : B_{k,\mathrm{rem}}^{(\ell)}>0\}.
\label{eq:downlink-payload-active}
\end{equation}
As users finish their payloads they leave the active set, which changes both
the interference structure and the feasible rate region of the remaining users.

\section{Uplink Formulation}

\subsection{Effective Uplink Signal Model}

In the uplink, user $k$ in block $\ell$ is represented through the effective
MIMO relation
\begin{equation}
\mathbf y_{k,\ell}^{(\UL)}
=
\Hc_{k,\ell}^{(\UL)}
\F_{k,\ell}^{(\UL)}
\mathbf s_{k,\ell}
+
\mathbf z_{k,\ell}^{(\UL)},
\label{eq:ul-signal}
\end{equation}
where
\begin{equation}
\Hc_{k,\ell}^{(\UL)} \in \C^{N_{r,k}\times N_{t,k}},
\qquad
\F_{k,\ell}^{(\UL)} \in \C^{N_{t,k}\times d_k},
\label{eq:ul-dimensions}
\end{equation}
$\mathbf s_{k,\ell}$ is the transmitted symbol vector with identity covariance,
and $\mathbf z_{k,\ell}^{(\UL)}$ is an AWGN vector. Hence its covariance matrix
is
\begin{equation}
\Sig_{k,\ell}^{(\UL)}
\sigma_k^2\I.
\label{eq:ul-noise-cov}
\end{equation}

The desired-signal covariance is
\begin{equation}
\Smat_{k,\ell}^{(\UL)}
=
\Hc_{k,\ell}^{(\UL)}
\F_{k,\ell}^{(\UL)}
\bigl(\F_{k,\ell}^{(\UL)}\bigr)^H
\bigl(\Hc_{k,\ell}^{(\UL)}\bigr)^H,
\label{eq:ul-signal-cov}
\end{equation}
and the whitened matrix becomes
\begin{equation}
\A_{k,\ell}^{(\UL)}
=
\bigl(\Sig_{k,\ell}^{(\UL)}\bigr)^{-1/2}
\Smat_{k,\ell}^{(\UL)}
\bigl(\Sig_{k,\ell}^{(\UL)}\bigr)^{-1/2}.
\label{eq:ul-A}
\end{equation}
Substituting \eqref{eq:ul-A} into
\eqref{eq:common-C}--\eqref{eq:common-rate} gives
$C_{k,\ell}^{(\UL)}$, $V_{k,\ell}^{(\UL)}$, and
$R_{k,\ell}^{(\UL)}$.

\subsection{Uplink Payload Formulation}

In payload communication, the number of blocks required by user $k$ is not
known in advance. Let $L_k$ denote the number of blocks eventually used by user
$k$. This terminal block index is induced by the resulting optimization and is
therefore an output of the optimization rather than an independent optimization
parameter. The uplink payload problem is
\begin{subequations}
\begin{align}
\min_{\{b_{k,\ell},n_{k,\ell},\F_{k,\ell}^{(\UL)}\}}
\quad
&
\sum_{k\in\Kset}\frac{1}{f_{s,k}}\sum_{\ell=1}^{L_k} n_{k,\ell}
\label{eq:ul-payload-obj}
\\
\text{s.t.}\quad
&
\frac{b_{k,\ell}}{n_{k,\ell}}
\le
R_{k,\ell}^{(\UL)},
\qquad \forall k,\ell,
\label{eq:ul-payload-fbl}
\\
&
\sum_{\ell=1}^{L_k} b_{k,\ell} = B_k,
\qquad \forall k,
\label{eq:ul-payload-conservation}
\\
&
n_{k,\ell} \in \{1,\dots,T_k\},
\qquad \forall k,\ell,
\label{eq:ul-payload-n}
\\
&
\norm{\F_{k,\ell}^{(\UL)}}_F^2 \le P_k,
\qquad \forall k,\ell.
\label{eq:ul-payload-power}
\end{align}
\end{subequations}

The remaining-bits recursion is
\begin{equation}
B_{k,\mathrm{rem}}^{(1)}=B_k,
\qquad
B_{k,\mathrm{rem}}^{(\ell+1)}=B_{k,\mathrm{rem}}^{(\ell)}-b_{k,\ell},
\label{eq:ul-remaining-recursion}
\end{equation}
and the scheduler terminates once $B_{k,\mathrm{rem}}^{(\ell)}=0$.

\section{Downlink Formulation}

\subsection{Interference-Coupled Downlink Signal Model}

In the downlink, all active users in block $\ell$ are transmitted
simultaneously. The received signal of user $k$ is
\begin{equation}
\mathbf y_{k,\ell}^{(\DL)}
=
\Hc_{k,\ell}^{(\DL)}\F_{k,\ell}^{(\DL)}\mathbf s_{k,\ell}
+
\sum_{j\in\Aset_\ell\setminus\{k\}}
\Hc_{k,\ell}^{(\DL)}\F_{j,\ell}^{(\DL)}\mathbf s_{j,\ell}
+
\mathbf z_{k,\ell}^{(\DL)},
\label{eq:dl-signal}
\end{equation}
where
\begin{equation}
\Hc_{k,\ell}^{(\DL)} \in \C^{N_{r,k}\times N_b},
\qquad
\F_{k,\ell}^{(\DL)} \in \C^{N_b\times d_k},
\label{eq:dl-dimensions}
\end{equation}
and $\mathbf z_{k,\ell}^{(\DL)}$ is receiver noise.

The desired-signal covariance of user $k$ is
\begin{equation}
\Smat_{k,\ell}^{(\DL)}
=
\Hc_{k,\ell}^{(\DL)}
\F_{k,\ell}^{(\DL)}
\bigl(\F_{k,\ell}^{(\DL)}\bigr)^H
\bigl(\Hc_{k,\ell}^{(\DL)}\bigr)^H.
\label{eq:dl-signal-cov}
\end{equation}
Its interference-plus-noise covariance is
\begin{equation}
\Sig_{k,\ell}^{(\DL)}
=
\sigma_k^2\I
+
\sum_{j\in\Aset_\ell\setminus\{k\}}
\Hc_{k,\ell}^{(\DL)}
\F_{j,\ell}^{(\DL)}
\bigl(\F_{j,\ell}^{(\DL)}\bigr)^H
\bigl(\Hc_{k,\ell}^{(\DL)}\bigr)^H.
\label{eq:dl-disturbance-cov}
\end{equation}

Hence the downlink whitened matrix is
\begin{equation}
\A_{k,\ell}^{(\DL)}
=
\bigl(\Sig_{k,\ell}^{(\DL)}\bigr)^{-1/2}
\Smat_{k,\ell}^{(\DL)}
\bigl(\Sig_{k,\ell}^{(\DL)}\bigr)^{-1/2},
\label{eq:dl-A}
\end{equation}
which leads directly to
$C_{k,\ell}^{(\DL)}$, $V_{k,\ell}^{(\DL)}$, and
$R_{k,\ell}^{(\DL)}$ through
\eqref{eq:common-C}--\eqref{eq:common-rate}.

\subsection{Downlink Payload Formulation}

In payload communication, the active-user set is induced by the remaining
payload:
\begin{equation}
\Aset_\ell
=
\left\{
k\in\Kset : B_{k,\mathrm{rem}}^{(\ell)}>0
\right\}.
\label{eq:dl-active-payload}
\end{equation}

The network-level problem is
\begin{subequations}
\begin{align}
\min_{\{b_{k,\ell},n_{k,\ell},\F_{k,\ell}^{(\DL)}\}}
\quad
&
\sum_{k\in\Kset}\frac{1}{f_{s,k}}\sum_{\ell} n_{k,\ell}
\label{eq:dl-payload-obj}
\\
\text{s.t.}\quad
&
\frac{b_{k,\ell}}{n_{k,\ell}}
\le
R_{k,\ell}^{(\DL)},
\qquad \forall k,\ell,
\label{eq:dl-payload-fbl}
\\
&
\sum_{\ell} b_{k,\ell} = B_k,
\qquad \forall k,
\label{eq:dl-payload-conservation}
\\
&
n_{k,\ell} \in \{1,\dots,T_k\},
\qquad \forall k,\ell,
\label{eq:dl-payload-n}
\\
&
\sum_{k\in\Aset_\ell}\norm{\F_{k,\ell}^{(\DL)}}_F^2 \le P_B,
\qquad \forall \ell.
\label{eq:dl-payload-power}
\end{align}
\end{subequations}

The same remaining-bits recursion as in
\eqref{eq:ul-remaining-recursion} applies, but now the block service of one
user depends on the simultaneous service decisions of the others.

\section{Training Only Method}

\subsection{Common Primal-Dual Structure}

The \emph{Training Only Method} performs optimization directly on the current
channel realization. There is no offline learning stage. Instead, each decision
block is solved online through a constrained primal-dual update, with
$b_{k,\ell}^{\mathrm{req}} = B_{k,\mathrm{rem}}^{(\ell)}$.

For uplink block $(k,\ell)$ with bits-to-send value
$b_{k,\ell}^{\mathrm{req}}$, define the rate constraint residual
\begin{equation}
g_{k,\ell}^{(\UL,r)}
=
\frac{b_{k,\ell}^{\mathrm{req}}}{n_{k,\ell}} - R_{k,\ell}^{(\UL)},
\label{eq:ul-rate-residual}
\end{equation}
and the power residual
\begin{equation}
g_{k,\ell}^{(\UL,p)}
=
\norm{\F_{k,\ell}^{(\UL)}}_F^2 - P_k.
\label{eq:ul-power-residual}
\end{equation}

For the downlink block, define
\begin{equation}
g_{k,\ell}^{(\DL,r)}
=
\frac{b_{k,\ell}^{\mathrm{req}}}{n_{k,\ell}} - R_{k,\ell}^{(\DL)},
\qquad k\in\Aset_\ell,
\label{eq:dl-rate-residual}
\end{equation}
and the common block-power residual
\begin{equation}
g_{\ell}^{(\DL,p)}
=
\sum_{k\in\Aset_\ell}\norm{\F_{k,\ell}^{(\DL)}}_F^2 - P_B.
\label{eq:dl-power-residual}
\end{equation}

The blockwise Lagrangian is then formed by combining a rate-maximizing primal
objective with nonnegative dual penalties. In uplink,
\begin{equation}
\mathcal{L}_{k,\ell}^{(\UL)}
=
-R_{k,\ell}^{(\UL)}
+
\lambda_{k,\ell}^{(\UL,r)}\bigl[g_{k,\ell}^{(\UL,r)}\bigr]_+
+
\lambda_{k,\ell}^{(\UL,p)}\bigl[g_{k,\ell}^{(\UL,p)}\bigr]_+,
\label{eq:ul-lagrangian}
\end{equation}
while in downlink,
\begin{equation}
\mathcal{L}_{\ell}^{(\DL)}
=
-\sum_{k\in\Aset_\ell} R_{k,\ell}^{(\DL)}
+
\sum_{k\in\Aset_\ell}
\lambda_{k,\ell}^{(\DL,r)}\bigl[g_{k,\ell}^{(\DL,r)}\bigr]_+
+
\lambda_{\ell}^{(\DL,p)}\bigl[g_{\ell}^{(\DL,p)}\bigr]_+.
\label{eq:dl-lagrangian}
\end{equation}

The primal variables are updated to reduce the Lagrangian, and the dual
variables are updated to penalize persistent infeasibility. To monitor the
solve, the method tracks three KKT residuals.

For uplink block $(k,\ell)$,
\begin{align}
r_{p,k,\ell}^{(\UL)}
&=
\max\!\left\{
\bigl[g_{k,\ell}^{(\UL,r)}\bigr]_+,\,
\bigl[g_{k,\ell}^{(\UL,p)}\bigr]_+
\right\},
\label{eq:ul-rp}
\\
r_{c,k,\ell}^{(\UL)}
&=
\max\!\left\{
\lambda_{k,\ell}^{(\UL,r)}\bigl[g_{k,\ell}^{(\UL,r)}\bigr]_+,\,
\lambda_{k,\ell}^{(\UL,p)}\bigl[g_{k,\ell}^{(\UL,p)}\bigr]_+
\right\},
\label{eq:ul-rc}
\\
r_{s,k,\ell}^{(\UL)}
&=
\frac{
\norm{
\F_{k,\ell}^{(\UL),(t)}-\F_{k,\ell}^{(\UL),(t-1)}
}_F
}{
\max\!\left\{
\norm{\F_{k,\ell}^{(\UL),(t-1)}}_F,\,
\delta_0
\right\}
},
\label{eq:ul-rs}
\end{align}
where $\delta_0>0$ is a small safeguard constant.

For downlink block $\ell$,
\begin{align}
r_{p,\ell}^{(\DL)}
&=
\max\!\left\{
\max_{k\in\Aset_\ell}\bigl[g_{k,\ell}^{(\DL,r)}\bigr]_+,\,
\bigl[g_{\ell}^{(\DL,p)}\bigr]_+
\right\},
\label{eq:dl-rp}
\\
r_{c,\ell}^{(\DL)}
&=
\max\!\left\{
\max_{k\in\Aset_\ell}
\lambda_{k,\ell}^{(\DL,r)}\bigl[g_{k,\ell}^{(\DL,r)}\bigr]_+,\,
\lambda_{\ell}^{(\DL,p)}\bigl[g_{\ell}^{(\DL,p)}\bigr]_+
\right\},
\label{eq:dl-rc}
\\
r_{s,\ell}^{(\DL)}
&=
\max_{k\in\Aset_\ell}
\frac{
\norm{
\F_{k,\ell}^{(\DL),(t)}-\F_{k,\ell}^{(\DL),(t-1)}
}_F
}{
\max\!\left\{
\norm{\F_{k,\ell}^{(\DL),(t-1)}}_F,\,
\delta_0
\right\}
},
\label{eq:dl-rs}
\end{align}

The online solve is declared converged when
\begin{equation}
r_p \le \epsilon_p,
\qquad
r_c \le \epsilon_c,
\qquad
r_s \le \epsilon_s,
\label{eq:kkt-convergence}
\end{equation}
and it is declared stationary infeasible when
\begin{equation}
r_s \le \epsilon_s
\qquad \text{but} \qquad
r_p > \epsilon_p.
\label{eq:kkt-stationary-infeasible}
\end{equation}

\subsection{Uplink Training Only Method}

The uplink training-only method is fundamentally sequential, but it follows the
same \emph{converge first, allocate second, reduce third} logic as the
downlink method. For each user-local block, the solver first converges the
full-sub-blocklength uplink precoder at $n_{k,\ell}=T_k$, then computes the
supported bits, and only then attempts sub-blocklength reduction if the entire
payload remainder has already been met. The
full-block service is
\begin{equation}
b_{k,\ell}^{(T)}
=
\min\!\left\{
B_{k,\mathrm{rem}}^{(\ell)},
\floor{T_k\,R_{k,\ell}^{(\UL)}(T_k)}
\right\}.
\label{eq:ul-full-block-service}
\end{equation}



Sub-blocklength reduction is attempted only after the current payload remainder
has been fully met, namely when
$b_{k,\ell}^{(T)} = B_{k,\mathrm{rem}}^{(\ell)}$. If the current payload remainder is
only partially served, the block is committed at $n_{k,\ell}=T_k$ and the
method moves to the next request without trying to reduce sub-blocklength.

\begin{algorithm}[H]
\caption{Training-Only Uplink User Solver}
\label{alg:ul-train-only}
\small
\begin{algorithmic}[1]
\Require User $k$, payload $B_k$, tolerances $\epsilon_p,\epsilon_c,\epsilon_s$
\Ensure Block sequence $\{(b_{k,\ell},n_{k,\ell},\F_{k,\ell}^{(\UL)})\}$
\State Set $B_{k,\mathrm{rem}}^{(1)} \gets B_k$
\State Initialize $\ell \gets 1$
\While{$B_{k,\mathrm{rem}}^{(\ell)} > 0$}
    \State Set $b_{k,\ell}^{\mathrm{req}} \gets B_{k,\mathrm{rem}}^{(\ell)}$
    \State Set $n_{k,\ell}\gets T_k$ and initialize the primal-dual state $(\F_{k,\ell}^{(\UL)},\lambda_{k,\ell}^{(\UL,r)},\lambda_{k,\ell}^{(\UL,p)})$
    \Repeat
        \State Update $\F_{k,\ell}^{(\UL)}$ using the uplink Lagrangian \eqref{eq:ul-lagrangian}
        \State Update $\lambda_{k,\ell}^{(\UL,r)}$ and $\lambda_{k,\ell}^{(\UL,p)}$
        \State Measure $r_{p,k,\ell}^{(\UL)}$, $r_{c,k,\ell}^{(\UL)}$, and $r_{s,k,\ell}^{(\UL)}$ from \eqref{eq:ul-rp}--\eqref{eq:ul-rs}
    \Until{the uplink solve satisfies \eqref{eq:kkt-convergence} or \eqref{eq:kkt-stationary-infeasible}}
    \State Restore the best accepted state: feasible if available, otherwise minimum-violation
    \State Compute $b_{k,\ell}^{(T)}=\min\!\left\{b_{k,\ell}^{\mathrm{req}},\floor{T_k R_{k,\ell}^{(\UL)}(T_k)}\right\}$
    \If{$b_{k,\ell}^{(T)} = 0$}
        \State Set $b_{k,\ell}\gets 0$ and keep $n_{k,\ell}=T_k$
        \State Declare the remaining payload infeasible under the current state and stop
    \ElsIf{$b_{k,\ell}^{(T)} < b_{k,\ell}^{\mathrm{req}}$}
        \State Commit $b_{k,\ell}=b_{k,\ell}^{(T)}$ and set $n_{k,\ell}=T_k$
    \Else
        \State Set $b_{k,\ell}=b_{k,\ell}^{\mathrm{req}}$
        \State Store the current accepted state
        \While{a smaller candidate sub-blocklength exists}
            \State Tentatively replace $n_{k,\ell}$ by the next smaller candidate
            \State Re-solve the candidate state with the same primal-dual loop above
            \State Restore the best candidate state: feasible if available, otherwise the previous accepted state
            \If{$\dfrac{b_{k,\ell}}{n_{k,\ell}} \le R_{k,\ell}^{(\UL)}$ and $\norm{\F_{k,\ell}^{(\UL)}}_F^2 \le P_k$}
                \State Accept the candidate and store the updated accepted state
            \Else
                \State Stop reduction and keep the previous accepted state
            \EndIf
        \EndWhile
        \State Commit the smallest accepted $n_{k,\ell}$ and corresponding precoder state
    \EndIf
    \State Update $B_{k,\mathrm{rem}}^{(\ell+1)} \gets B_{k,\mathrm{rem}}^{(\ell)} - b_{k,\ell}$
    \State $\ell \gets \ell + 1$
\EndWhile
\end{algorithmic}
\end{algorithm}

Algorithm~\ref{alg:ul-train-only} solves the uplink payload problem by
updating the remaining workload after every committed block.

\subsection{Downlink Training Only Method}

In the downlink, a whole active-user block must be solved jointly. The same
three-phase logic is applied, but now on a coupled active-user set. The method
first converges the joint full-sub-blocklength precoder at the coherence limits
$\{T_k\}_{k\in\Aset_\ell}$, then computes the supported bits of all requested
users, removes users with zero supported service, and only then reduces
sub-blocklength for the users whose full payload remainders are already
supported. This produces the full-block supported service
\begin{equation}
b_{k,\ell}^{(T)}
=
\min\!\left\{
b_{k,\ell}^{\mathrm{req}},
\floor{T_k\,R_{k,\ell}^{(\DL)}(T_k)}
\right\},
\qquad k\in\Aset_\ell.
\label{eq:dl-full-block-service}
\end{equation}

Users with $b_{k,\ell}^{(T)}=0$ are removed from transmission in that block.
For users that satisfy
$b_{k,\ell}^{(T)}=b_{k,\ell}^{\mathrm{req}}$, the method tries to reduce
sub-blocklength while keeping the already committed bits fixed. Since the block
is coupled, a candidate reduction may invalidate another user's committed
service. The method therefore re-optimizes the affected active-user subset
whenever necessary and accepts a candidate reduction only if all committed bits
remain feasible.

\begin{algorithm}[H]
\caption{Training-Only Downlink Block Solver}
\label{alg:dl-train-only}
\small
\begin{algorithmic}[1]
\Require Active-user set $\Aset_\ell$, remaining payloads $\{B_{k,\mathrm{rem}}^{(\ell)}\}_{k\in\Aset_\ell}$, tolerances $\epsilon_p,\epsilon_c,\epsilon_s$
\Ensure Block plan $\{(b_{k,\ell},n_{k,\ell},\F_{k,\ell}^{(\DL)})\}_{k\in\Aset_\ell}$
\State Store the requested-user set as $\Aset_\ell^{\mathrm{req}}\gets \Aset_\ell$
\State Set the current transmit set $\Aset_\ell^{\mathrm{tx}}\gets \Aset_\ell^{\mathrm{req}}$
\For{each $k\in\Aset_\ell^{\mathrm{req}}$}
    \State Set $b_{k,\ell}^{\mathrm{req}} \gets B_{k,\mathrm{rem}}^{(\ell)}$
    \State Set $n_{k,\ell}\gets T_k$
\EndFor
\State Initialize the active-user precoders and multipliers on $\Aset_\ell^{\mathrm{tx}}$
\Repeat
    \State Update the active-user precoders using the block Lagrangian \eqref{eq:dl-lagrangian}
    \State Update the multipliers $\{\lambda_{k,\ell}^{(\DL,r)}\}$ and $\lambda_{\ell}^{(\DL,p)}$
    \State Measure $r_{p,\ell}^{(\DL)}$, $r_{c,\ell}^{(\DL)}$, and $r_{s,\ell}^{(\DL)}$ from \eqref{eq:dl-rp}--\eqref{eq:dl-rs}
\Until{the downlink solve satisfies \eqref{eq:kkt-convergence} or \eqref{eq:kkt-stationary-infeasible}}
\State Restore the best accepted block state: feasible if available, otherwise minimum-violation
\Repeat
    \For{each requested user $k\in\Aset_\ell^{\mathrm{req}}$}
        \If{$k\in\Aset_\ell^{\mathrm{tx}}$}
            \State Set $b_{k,\ell}\gets \min\!\left\{b_{k,\ell}^{\mathrm{req}},\floor{n_{k,\ell}R_{k,\ell}^{(\DL)}(\mathbf n_\ell)}\right\}$
        \Else
            \State Set $b_{k,\ell}\gets 0$
        \EndIf
    \EndFor
    \If{every user in $\Aset_\ell^{\mathrm{tx}}$ satisfies $b_{k,\ell}>0$}
        \State Stop pruning zero-service users
    \Else
        \State Remove from $\Aset_\ell^{\mathrm{tx}}$ every user with $b_{k,\ell}=0$
        \If{$\Aset_\ell^{\mathrm{tx}}\neq\varnothing$}
            \State Re-solve the reduced active-user block with the same primal-dual loop above
            \State Restore the best accepted reduced-set state: feasible if available, otherwise minimum-violation
        \EndIf
    \EndIf
\Until{every remaining active user receives positive full-state service or $\Aset_\ell^{\mathrm{tx}}=\varnothing$}
\For{each user $k\in\Aset_\ell^{\mathrm{tx}}$ satisfying $b_{k,\ell}=b_{k,\ell}^{\mathrm{req}}$}
    \State Store the current accepted block state
    \While{a smaller candidate sub-blocklength exists for user $k$}
        \State Tentatively decrease $n_{k,\ell}$ while holding the other active-user sub-blocklengths fixed
        \State Re-optimize the affected active-user block with the same primal-dual loop above
        \State Restore the best candidate block state: feasible if available, otherwise the previous accepted state
        \If{$\dfrac{b_{j,\ell}}{n_{j,\ell}} \le R_{j,\ell}^{(\DL)}$ for all $j\in\Aset_\ell^{\mathrm{tx}}$ and $\sum_{j\in\Aset_\ell^{\mathrm{tx}}}\norm{\F_{j,\ell}^{(\DL)}}_F^2 \le P_B$}
            \State Accept the candidate and store the updated accepted state
        \Else
            \State Stop reducing user $k$
        \EndIf
    \EndWhile
\EndFor
\For{each requested user $k\in\Aset_\ell^{\mathrm{req}}$}
    \State Update $B_{k,\mathrm{rem}}^{(\ell+1)} \gets B_{k,\mathrm{rem}}^{(\ell)} - b_{k,\ell}$
\EndFor
\State Commit the final joint block plan
\end{algorithmic}
\end{algorithm}

The training-only downlink method is therefore a \emph{converge first, allocate
second, reduce third} procedure. This ordering is essential because the
finite-sub-blocklength rate of each user depends on the full active-user beam set.

\section{Training + Testing Method}

\subsection{Method Principle}

The \emph{Training + Testing Method} separates optimization into two stages.
During training, a precoder mapping is learned from a collection of queries
produced by the same finite-sub-blocklength scheduling logic that will later be
used at test time. During testing, the online optimization step is replaced by
model inference, followed by the same feasibility checks and the same
sub-blocklength-reduction logic.

The central training object is a query. In general form, a query
$q\in\Qset$ contains
\begin{equation}
\begin{aligned}
q
=
\bigl(
&
\text{channel state},\,
\text{active-user state},\,
\text{candidate sub-blocklength state},\\
&
\text{bits-to-send state},\,
\text{covariance state},\,
\text{reliability target}
\bigr).
\end{aligned}
\label{eq:general-query}
\end{equation}

Thus, a query is simply one scheduling state at which the model is asked to
predict a feasible finite-sub-blocklength transmission strategy. In the present
formulation, the training loss uses \emph{uniform-per-query} averaging, so every
stored query contributes equally to the empirical objective.

\subsection{Uplink Training + Testing Method}

\subsubsection{Training stage}

For uplink user $k$, a query at block $\ell$ has the form
\begin{equation}
q_{k,\ell}^{(\UL)}
=
\bigl(
\Hc_{k,\ell}^{(\UL)},
\Sig_{k,\ell}^{(\UL)},
n,
b_{k,\ell}^{\mathrm{req}},
\varepsilon_k
\bigr),
\label{eq:ul-query}
\end{equation}
where $n$ is the candidate sub-blocklength being queried and
$b_{k,\ell}^{\mathrm{req}}$ is the number of bits to send in that state. The
query set $\Qset_k^{(\UL)}$ is built by running the same causal block logic
used by the scheduler itself and recording every visited candidate state. In
payload communication, the number of bits to send is
$b_{k,\ell}^{\mathrm{req}} = B_{k,\mathrm{rem}}^{(\ell)}$.

Let $\theta_k$ denote the trainable parameters of the uplink model for user $k$.
The uniform-per-query training objective is
\begin{equation}
\mathcal{J}_k^{(\UL)}(\theta_k)
=
\frac{1}{\lvert\Qset_k^{(\UL)}\rvert}
\sum_{q\in\Qset_k^{(\UL)}}
\left(
-R(q;\theta_k)
+
\lambda_k^{(r)}\Bigl[\frac{b^{\mathrm{req}}(q)}{n(q)}-R(q;\theta_k)\Bigr]_+
+
\lambda_k^{(p)}\bigl[P(q,\theta_k)-P_k\bigr]_+
\right),
\label{eq:ul-training-objective}
\end{equation}
where $R(q;\theta_k)$ is the finite-sub-blocklength rate induced by the predicted
uplink precoder, $n(q)$ is the queried sub-blocklength, $b^{\mathrm{req}}(q)$ is
the bits-to-send value encoded in query $q$, and $P(q,\theta_k)$ is the
transmit power induced by the predicted precoder. The coefficient
$\lambda_k^{(r)}$ penalizes rate shortfall relative to the bits-to-send value,
whereas $\lambda_k^{(p)}$ penalizes power-budget violation.

\begin{algorithm}[H]
\caption{Uplink Training Stage}
\label{alg:ul-train-stage}
\small
\begin{algorithmic}[1]
\Require Training episodes for uplink user $k$
\Ensure Query set $\Qset_k^{(\UL)}$ and trained parameters $\theta_k$
\State Initialize $\Qset_k^{(\UL)}\gets\varnothing$
\For{each training episode}
    \State Initialize the episode remainder state from the episode payload
    \For{each visited block $\ell$ of user $k$}
        \State Set $b_{k,\ell}^{\mathrm{req}} \gets B_{k,\mathrm{rem}}^{(\ell)}$
        \State Set the first candidate sub-blocklength to $n\gets T_k$
        \While{$n$ is under consideration}
            \State Form the query $q_{k,\ell}^{(\UL)}$ in \eqref{eq:ul-query} and add it to $\Qset_k^{(\UL)}$
            \State Evaluate the current state and compute
            \[
            b_{k,\ell}^{(n)}
            =
            \min\!\left\{
            b_{k,\ell}^{\mathrm{req}},
            \floor{n\,R_{k,\ell}^{(\UL)}(n)}
            \right\}
            \]
            \If{$b_{k,\ell}^{(n)} < b_{k,\ell}^{\mathrm{req}}$}
                \State Stop collecting additional queries for this sub-block
            \Else
                \State Replace $n$ by the next smaller candidate sub-blocklength
            \EndIf
        \EndWhile
        \State Advance the episode remainder by the committed service
    \EndFor
\EndFor
\State Train $\theta_k$ by minimizing \eqref{eq:ul-training-objective}
\end{algorithmic}
\end{algorithm}

\subsubsection{Testing stage}

At test time, the trained model replaces the online uplink solve. The block is
first examined at the full available sub-blocklength. If the full block does not
support the complete request, the scheduler commits only the supported service
at $n=T_k$. If the full block already supports the entire request, the
sub-blocklength is reduced and the model is queried again at smaller candidate
sub-blocklengths until the smallest feasible one is identified.

\begin{algorithm}[H]
\caption{Uplink Testing Stage}
\label{alg:ul-test-stage}
\small
\begin{algorithmic}[1]
\Require Trained uplink model $\theta_k$, user $k$, payload $B_k$
\Ensure Final uplink block plan
\State Set $B_{k,\mathrm{rem}}^{(1)} \gets B_k$
\For{each visited block $\ell$}
    \State Set $b_{k,\ell}^{\mathrm{req}} \gets B_{k,\mathrm{rem}}^{(\ell)}$
    \State Infer the uplink precoder at $n=T_k$
    \State Compute the full-block supported service
    \[
    b_{k,\ell}^{(T)}
    =
    \min\!\left\{
    b_{k,\ell}^{\mathrm{req}},
    \floor{T_k\,R_{k,\ell}^{(\UL)}(T_k)}
    \right\}
    \]
    \If{$b_{k,\ell}^{(T)} < b_{k,\ell}^{\mathrm{req}}$}
        \State Commit $n_{k,\ell}=T_k$ and $b_{k,\ell}=b_{k,\ell}^{(T)}$
    \Else
        \State Set $b_{k,\ell}=b_{k,\ell}^{\mathrm{req}}$
        \State Set $n_{k,\ell}\gets T_k$
        \State Store the current accepted state
        \While{a smaller candidate sub-blocklength exists}
            \State Tentatively replace $n_{k,\ell}$ by the next smaller candidate
            \State Infer the uplink precoder for the updated query
            \If{$\dfrac{b_{k,\ell}}{n_{k,\ell}} \le R_{k,\ell}^{(\UL)}(n_{k,\ell})$ and $P(q_{k,\ell}^{(\UL)},\theta_k)\le P_k$}
                \State Accept the reduction and store the updated accepted state
            \Else
                \State Stop reduction and keep the previous accepted state
            \EndIf
        \EndWhile
        \State Commit the smallest accepted $n_{k,\ell}$ and corresponding precoder
    \EndIf
    \State Update $B_{k,\mathrm{rem}}^{(\ell+1)} \gets B_{k,\mathrm{rem}}^{(\ell)} - b_{k,\ell}$
\EndFor
\end{algorithmic}
\end{algorithm}

\subsection{Downlink Training + Testing Method}

\subsubsection{Training stage}

The downlink training state is joint across the active users of a block. A
query for block $\ell$ can be written as
\begin{equation}
q_{\ell}^{(\DL)}
=
\bigl(
\{\Hc_{k,\ell}^{(\DL)}\}_{k\in\Aset_\ell},
\Aset_\ell,
\{n_{k,\ell}\}_{k\in\Aset_\ell},
\{b_{k,\ell}^{\mathrm{req}}\}_{k\in\Aset_\ell},
\{\Sig_{k,\ell}^{(\DL)}\}_{k\in\Aset_\ell},
\{\varepsilon_k\}_{k\in\Aset_\ell}
\bigr),
\label{eq:dl-query}
\end{equation}
where $\Aset_\ell$ is the active-user set of the block. The query set
$\Qset^{(\DL)}$ is built causally by running the same joint scheduling logic
used at test time and recording every visited joint state. In payload
communication, the number of bits to send is
$b_{k,\ell}^{\mathrm{req}} = B_{k,\mathrm{rem}}^{(\ell)}$.

Let $\theta=\{\theta_k\}_{k=1}^K$ denote the trainable parameters of the
downlink model collection. The uniform-per-query training objective is
\begin{equation}
\begin{aligned}
\mathcal{J}^{(\DL)}(\theta)
=
\frac{1}{\lvert\Qset^{(\DL)}\rvert}
\sum_{q\in\Qset^{(\DL)}}
\Biggl(
&-\sum_{k\in\Aset(q)} R_k(q;\theta)\\
&+ \sum_{k\in\Aset(q)}
\lambda_k^{(r)}
\Bigl[\frac{b_k^{\mathrm{req}}(q)}{n_k(q)}-R_k(q;\theta)\Bigr]_+\\
&+ \lambda^{(p)}
\Bigl[
\sum_{k\in\Aset(q)}P_k(q,\theta)-P_B
\Bigr]_+
\Biggr)
\Biggr).
\end{aligned}
\label{eq:dl-training-objective}
\end{equation}
where $\Aset(q)$ is the active-user set encoded in query $q$, $n_k(q)$ is the
sub-blocklength associated with active user $k$ in that query,
$b_k^{\mathrm{req}}(q)$ is the number of bits to send for active user $k$, and
$P_k(q,\theta)$ is the transmit power induced for active user $k$ by the
predicted downlink precoder. The penalty coefficients $\lambda_k^{(r)}$ and
$\lambda^{(p)}$ enforce the rate constraints and the total base-station power
constraint, respectively.

\begin{algorithm}[H]
\caption{Downlink Training Stage}
\label{alg:dl-train-stage}
\small
\begin{algorithmic}[1]
\Require Training episodes
\Ensure Query set $\Qset^{(\DL)}$ and trained parameters $\theta$
\State Initialize $\Qset^{(\DL)}\gets\varnothing$
\For{each training episode}
    \For{each visited communication block $\ell$}
        \State Determine the requested-user set $\Aset_\ell^{\mathrm{req}}$
        \State Set the current transmit set $\Aset_\ell^{\mathrm{tx}}\gets \Aset_\ell^{\mathrm{req}}$
        \For{each $k\in\Aset_\ell^{\mathrm{req}}$}
            \State Set $b_{k,\ell}^{\mathrm{req}} \gets B_{k,\mathrm{rem}}^{(\ell)}$
            \State Initialize $n_{k,\ell}\gets T_k$
        \EndFor
        \While{the current joint sub-blocklength state is under consideration}
            \State Form the joint query $q_{\ell}^{(\DL)}$ in \eqref{eq:dl-query} and add it to $\Qset^{(\DL)}$
            \For{each $k\in\Aset_\ell^{\mathrm{req}}$}
                \If{$k\in\Aset_\ell^{\mathrm{tx}}$}
                    \State Set $b_{k,\ell}\gets \min\!\left\{b_{k,\ell}^{\mathrm{req}},\floor{n_{k,\ell}R_{k,\ell}^{(\DL)}(\mathbf n_\ell)}\right\}$
                \Else
                    \State Set $b_{k,\ell}\gets 0$
                \EndIf
            \EndFor
            \If{some user in $\Aset_\ell^{\mathrm{tx}}$ satisfies $b_{k,\ell}=0$}
                \State Remove from $\Aset_\ell^{\mathrm{tx}}$ every user with $b_{k,\ell}=0$
            \ElsIf{every user in $\Aset_\ell^{\mathrm{tx}}$ satisfies $b_{k,\ell}<b_{k,\ell}^{\mathrm{req}}$}
                \State Stop collecting additional queries for this block
            \Else
                \State Decrease the candidate sub-blocklengths of users satisfying $b_{k,\ell}=b_{k,\ell}^{\mathrm{req}}$
            \EndIf
        \EndWhile
        \State Advance the episode remainders by the committed service
    \EndFor
\EndFor
\State Train $\theta$ by minimizing \eqref{eq:dl-training-objective}
\end{algorithmic}
\end{algorithm}

\subsubsection{Testing stage}

At test time, the learned downlink model is queried on the current active-user
state of the block. The full joint sub-blocklength vector is evaluated first. Users
with zero supported service are removed from the transmit set, and joint
sub-blocklength reduction is attempted only for users whose full bits-to-send
value is already supported at the current state.

\begin{algorithm}[H]
\caption{Downlink Testing Stage}
\label{alg:dl-test-stage}
\small
\begin{algorithmic}[1]
\Require Trained downlink model $\theta$, remaining payloads $\{B_{k,\mathrm{rem}}^{(\ell)}\}$
\Ensure Final downlink block plan
\For{each communication block $\ell$}
    \State Determine the requested-user set $\Aset_\ell^{\mathrm{req}}$
    \State Set the current transmit set $\Aset_\ell^{\mathrm{tx}}\gets \Aset_\ell^{\mathrm{req}}$
    \For{each $k\in\Aset_\ell^{\mathrm{req}}$}
        \State Set $b_{k,\ell}^{\mathrm{req}} \gets B_{k,\mathrm{rem}}^{(\ell)}$
        \State Initialize $n_{k,\ell}\gets T_k$
    \EndFor
    \State Infer the joint downlink precoders at the full sub-blocklength vector on $\Aset_\ell^{\mathrm{tx}}$
    \Repeat
        \For{each $k\in\Aset_\ell^{\mathrm{req}}$}
            \If{$k\in\Aset_\ell^{\mathrm{tx}}$}
                \State Set $b_{k,\ell}\gets \min\!\left\{b_{k,\ell}^{\mathrm{req}},\floor{n_{k,\ell}R_{k,\ell}^{(\DL)}(\mathbf n_\ell)}\right\}$
            \Else
                \State Set $b_{k,\ell}\gets 0$
            \EndIf
        \EndFor
        \If{some user in $\Aset_\ell^{\mathrm{tx}}$ satisfies $b_{k,\ell}=0$}
            \State Remove from $\Aset_\ell^{\mathrm{tx}}$ every user with $b_{k,\ell}=0$
            \If{$\Aset_\ell^{\mathrm{tx}}\neq\varnothing$}
                \State Re-infer the joint downlink precoders on the reduced transmit set
            \EndIf
        \Else
            \State Stop pruning zero-service users
        \EndIf
    \Until{every remaining active user receives positive full-state service or $\Aset_\ell^{\mathrm{tx}}=\varnothing$}
    \For{each user $k\in\Aset_\ell^{\mathrm{tx}}$ satisfying $b_{k,\ell}=b_{k,\ell}^{\mathrm{req}}$}
        \State Store the current accepted state
        \While{a smaller candidate sub-blocklength exists for that user}
            \State Tentatively decrease $n_{k,\ell}$ while holding the other active-user sub-blocklengths fixed
            \State Infer the joint downlink precoders for the updated state
            \If{$\dfrac{b_{j,\ell}}{n_{j,\ell}} \le R_{j,\ell}^{(\DL)}(\mathbf n_\ell)$ for all $j\in\Aset_\ell^{\mathrm{tx}}$ and $\sum_{j\in\Aset_\ell^{\mathrm{tx}}}P_j(q_{\ell}^{(\DL)},\theta)\le P_B$}
                \State Accept the reduction and store the updated accepted state
            \Else
                \State Stop reducing user $k$ and keep the previous accepted state
            \EndIf
        \EndWhile
    \EndFor
    \For{each requested user $k\in\Aset_\ell^{\mathrm{req}}$}
        \State Update $B_{k,\mathrm{rem}}^{(\ell+1)} \gets B_{k,\mathrm{rem}}^{(\ell)} - b_{k,\ell}$
    \EndFor
    \State Commit the final active-user set, sub-blocklengths, and transmitted bits of the block
\EndFor
\end{algorithmic}
\end{algorithm}

\section{Conclusion}

The report has recast the current work into a theory-centered structure. The
development proceeds from a common finite-sub-blocklength model and payload
communication setting, to separate uplink and downlink formulations, and then
to two solution frameworks.

The uplink problem is characterized by user-wise sequential block creation under
finite-sub-blocklength feasibility. The downlink problem is characterized by a
jointly coupled interference structure that requires block-level coordination
before blockwise service can be committed.

The \emph{Training Only Method} solves each decision state online through
constrained primal-dual optimization and feasibility-guided sub-blocklength
reduction. The \emph{Training + Testing Method} replaces online optimization by
learned precoder inference, but retains the same finite-sub-blocklength service
logic through query-based offline training, feasibility checks, and causal
sub-blocklength reduction.

Taken together, these formulations provide a complete theoretical description of
the current uplink and downlink methodology without relying on implementation
details or numerical experiment summaries.

\begin{thebibliography}{1}

\bibitem{polyanskiy2010}
Y.~Polyanskiy, H.~V.~Poor, and S.~Verd\'u,
``Channel coding rate in the finite sub-blocklength regime,''
\emph{IEEE Transactions on Information Theory},
vol.~56, no.~5, pp.~2307--2359, 2010.

\end{thebibliography}

\end{document}
